
\end{proof}
\end{proposition}
\begin{remark}
    Note that different orderings of indices of $\EE(\Ab\outprod\Ab)$ lead to the following tensor networks:  
    \begin{align*}
            \EE(\Ab\outprod\Ab)_{i_1i_2i_4i_3}
            &= 
            \EE(\Ab_{i_1i_2}\Ab_{i_4i_3}) 
    = \delta_{i_1i_4}\delta_{i_2i_3}\\
    &=
    \begin{cases}
        1 & \text{ if } i_1=i_4\text{ and } i_2=i_3\\
        0 & \text{ otherwise, }
    \end{cases} 
    =
    \begin{tikzpicture}[baseline=\baseline]
    \tikzset{tensor/.style = {minimum size = 0.5cm,shape = circle,thick,draw=black,inner sep = 0pt}, edge/.style = {thick,line width=.4mm},every loop/.style={}}
    \def\y{0}
    \def\x{0}
    \draw[edge] (\x,\y-0.1) to [out=100,in=80] (\x+2,\y-0.1);
    \draw[edge] (\x+0.15,\y-0.1) to [out=100,in=80] (\x+1.85,\y-0.1);
    \node[draw=none] () at (\x,\y-0.25) {\scalebox{0.7}{\indcolor{$i_1$}}};
    \node[draw=none] () at (\x+2.15,\y-0.25) {\scalebox{0.7}{\indcolor{$i_4$}}};
    \node[draw=none] () at (\x+0.35,\y-0.25) {\scalebox{0.7}{\indcolor{$i_2$}}};
    \node[draw=none] () at (\x+1.7,\y-0.25) {\scalebox{0.7}{\indcolor{$i_3$}}};
\end{tikzpicture},
    \end{align*}
and
\begin{align*}
\EE(\Ab\outprod\Ab)_{i_1i_4i_2i_3} 
&= 
\EE(\Ab_{i_1i_4}\Ab_{i_2i_3})= \delta_{i_1i_2}\delta_{i_4i_3}\\
&=
\begin{cases}
        1 & \text{ if } i_1=i_2\text{ and } i_3=i_4\\
        0 & \text{ otherwise, }
\end{cases}
=
\begin{tikzpicture}[baseline=\baseline]
    \tikzset{tensor/.style = {minimum size = 0.5cm,shape = circle,thick,draw=black,inner sep = 0pt}, edge/.style = {thick,line width=.4mm},every loop/.style={}}
    \def\y{0}
    \def\x{0}
    \draw[edge] (\x+0.15,\y-0.1) to [out=100,in=80] (\x+1.85,\y-0.1);
    \draw[edge] (\x+2,\y-0.1) to [out=100,in=80] (\x+3.65,\y-0.1);
    \node[draw=none] () at (\x+0.15,\y-0.25) {\scalebox{0.7}{\indcolor{$i_1$}}};
    \node[draw=none] () at (\x+1.85,\y-0.25) {\scalebox{0.7}{\indcolor{$i_2$}}};
    \node[draw=none] () at (\x+2.15,\y-0.25) {\scalebox{0.7}{\indcolor{$i_3$}}};
    \node[draw=none] () at (\x+3.65,\y-0.25) {\scalebox{0.7}{\indcolor{$i_4$}}};
\end{tikzpicture}.
\end{align*}
\end{remark}
Next, we show how Propositions~\ref{prop:inner-expectation} and~\ref{prop:outer-prod} can be used to derive a very simple tensor network derivation of the expected value of the product of a random matrix with Gaussian entries with itself. 
\begin{proposition}\label{prop:wishart-mean}
If $\Ab\in\RR^{m\times n}$ is a random matrix whose elements are drawn i.i.d from the standard normal distribution, then $\EE(\Ab^\ts\Ab) = m\Ib_n$. 
\begin{proof}
        Using Propositions~\ref{prop:inner-expectation} and~\ref{prop:outer-prod} we can write 
\begin{align*}
\EE(\Ab^\ts\Ab) 
&= 
\EE\left(\begin{tikzpicture}[baseline=\baseline]
    \tikzset{tensor/.style = {minimum size = 0.5cm,shape = circle,thick,draw=black,inner sep = 0pt}, edge/.style = {thick,line width=.4mm},every loop/.style={}}
    \def\y{0}
    \def\x{0}
    \node[tensor,fill=red!30!blue!30!white] (A1) at (\x,\y){\scalebox{0.85}{$\Ab$}};
    \node[tensor,fill=red!30!blue!30!white] (A2) at (\x+1,\y){\scalebox{0.85}{$\Ab$}};
    \draw[edge](A1) -- (A2) node [midway, above]{\scalebox{0.7}{\dimcolor{$m$}}};
    \draw[edge](A1) -- (\x-0.85,\y) node [midway, above]{\scalebox{0.7}{\dimcolor{$n$}}};
    \draw[edge](A2) -- (\x+1.85,\y) node [midway, above]{\scalebox{0.7}{\dimcolor{$n$}}};
\end{tikzpicture}\right)
=
\EE\left(\begin{tikzpicture}[baseline=\baseline]
    \tikzset{tensor/.style = {minimum size = 0.5cm,shape = circle,thick,draw=black,inner sep = 0pt}, edge/.style = {thick,line width=.4mm},every loop/.style={}}
    \def\y{0}
    \def\x{0}
    \node[tensor,fill=red!30!blue!30!white] (A2) at (\x,\y-0.5){\scalebox{0.85}{$\Ab$}};
    \node[tensor,fill=red!30!blue!30!white] (A1) at (\x,\y+0.5){\scalebox{0.85}{$\Ab$}};
    \draw[edge](\x+0.25,\y+0.4) -- (\x+1.25,\y+0.4) -- (\x+1.25,\y-0.4) -- (\x+0.25,\y-0.4);
    \draw[edge](\x+0.2,\y+0.65)--(\x+1.25,\y+0.65)node [midway, above]{\scalebox{0.7}{\dimcolor{$n$}}};
    \draw[edge](\x+0.2,\y-0.65)--(\x+1.25,\y-0.65)node [midway, below]{\scalebox{0.7}{\dimcolor{$n$}}};
    \node[draw=none] () at (\x+1.45,\y) {\scalebox{0.7}{\dimcolor{$m$}}};
    \draw[dotted](\x+0.55,\y+1)--(\x+0.55,\y-1);
\end{tikzpicture}\right)\\
\ \
&\textequal{Prop.~\ref{prop:inner-expectation}}
\ \
\begin{tikzpicture}[baseline=\baseline]
   \tikzset{tensor/.style = {minimum size = 0.5cm,shape = circle,thick,draw=black,inner sep = 0pt}, edge/.style = {thick,line width=.4mm},every loop/.style={}}
    \def\y{0}
    \def\x{0}
    \node[tensor,fill=red!30!blue!30!white] (A1) at (\x,\y+0.5){\scalebox{0.85}{$\Ab$}};
    \node[tensor,fill=red!30!blue!30!white] (A2) at (\x,\y-0.5){\scalebox{0.85}{$\Ab$}};
    \draw[edge,draw=blue](\x+0.25,\y+0.4) -- (\x+1.65,\y+0.4);
    \draw[edge, draw=red](\x+1.65,\y+0.4)-- (\x+1.95,\y+0.4) -- (\x+1.95,\y-0.4)--(\x+1.65,\y+0-0.4);
    \draw[edge, draw=blue](\x+1.65,\y-0.4) -- (\x+0.25,\y-0.4);
    \draw[edge, draw=blue](\x+0.2,\y+0.65)--(\x+1.95,\y+0.65)node [midway, above]{\scalebox{0.7}{\dimcolor{$n$}}};
    \draw[edge,draw=blue](\x+0.2,\y-0.65)--(\x+1.95,\y-0.65)node [midway, below]{\scalebox{0.7}{\dimcolor{$n$}}};
    \node[draw=none] () at (\x+2.15,\y) {\scalebox{0.7}{\dimcolor{$m$}}};
    \node at (\x-0.55,\y) {$\EE\left( \vphantom{\rule{0pt}{1cm}} \right.$};
    \node at (\x+0.55,\y) {$\left. \vphantom{\rule{0pt}{1cm}} \right)$};
    \node at (\x+1.15,\y) {$\EE\left( \vphantom{\rule{0pt}{1cm}} \right.$};
    \node at (\x+2.35,\y) {$\left. \vphantom{\rule{0pt}{1cm}} \right)$};
\end{tikzpicture}
~~\textequal{Prop.~\ref{prop:outer-prod}}
\ \ 
\begin{tikzpicture}[baseline=\baseline]
   \tikzset{tensor/.style = {minimum size = 0.5cm,shape = circle,thick,draw=black,inner sep = 0pt}, edge/.style = {thick,line width=.4mm},every loop/.style={}}
    \def\y{0}
    \def\x{0}
    \draw[edge,draw=blue](\x+1.35,\y+0.4) -- (\x+1.65,\y+0.4);
    \draw[edge, draw=red](\x+1.65,\y+0.4)-- (\x+1.95,\y+0.4) -- (\x+1.95,\y-0.4)--(\x+1.65,\y+0-0.4);
    \draw[edge, draw=blue](\x+1.35,\y-0.4) -- (\x+1.65,\y-0.4);
    \draw[edge, draw=blue](\x+1.35,\y-0.4) -- (\x+1.35,\y-0.4);
    \draw[edge, draw=blue](\x+1.35,\y+0.65)--(\x+1.95,\y+0.65);
    \draw[edge,draw=blue](\x+1.35,\y-0.65)--(\x+1.95,\y-0.65);
    \node[draw=none] () at (\x+2.15,\y) {\scalebox{0.7}{\dimcolor{$m$}}};
    \node[draw=none] () at (\x+1.35,\y) {\scalebox{0.7}{\dimcolor{$m$}}};
    \node[draw=none] () at (\x+0.85,\y) {\scalebox{0.7}{\dimcolor{$n$}}};
    \draw[edge, draw=blue](\x+1.35,\y+0.65)to[out=170,in=190](\x+1.35,\y-0.65);
    \draw[edge, draw=blue](\x+1.35,\y+0.4)to[out=190,in=170](\x+1.35,\y-0.4);
\end{tikzpicture}
= m\Ib_n.
\end{align*}
The final diagram is obtained by contracting the legs of size $m$. As observed, once the contraction is performed, the resulting expression is simply the trace of the identity matrix, which corresponds to the circle of size $m$ in the above equality.
\end{proof}
\end{proposition}
The previous proposition can be interpreted as a statement on the mean of a Wishart distribution.  
Recall that  sampling a matrix $\mat X$ from the Wishart distribution $W_p(\mat \Sigma, m)$, where $\mat \Sigma \in \RR^{n\times n}$ is a covariance matrix, is done by sampling each column of a matrix $\mat A\in\mathbb R^{m\times n}$ from the multivariate normal $\mathcal N(\vec 0 , \mat \Sigma)$ independently and setting $\mat X = \mat A^\top\mat A$. Hence, the matrix $\Ab ^\top \Ab$ in Proposition~\ref{prop:inner-expectation} follows the Wishart distribution  $W_p(\mat I,m)$.
\begin{remark} As a consequence of Proposition~\ref{prop:inner-expectation}, if $\Ab \in \RR^{m \times n}$ is a random matrix with entries drawn from the standard normal distribution, then $\EE(\norm{\Ab}_F^2) = mn$. 
This result can be extended to higher-order tensors whose entries are drawn i.i.d. from the standard normal distribution, e.g., if $\Tt\in\RR^{d_1\times d_2\times\cdots\times d_N}$ then $\EE(\norm{\Tt}_F^2) = d_1d_2\cdots d_N$.
\end{remark}
By leveraging the properties of random Gaussian matrices established in the previous propositions, we can compute the expected norm of their product. Once again, the proof is remarkably simple when expressed through tensor networks, whereas an explicit treatment using indices and summations would be considerably more intricate.
\begin{proposition}\label{prop:expected-norm}
    Let $\Ab\in\RR^{m\times r}$ and $\Bb\in\RR^{r\times n}$ be random matrices whose entries are independently and identically distributed (i.i.d.) according to the standard normal distribution with mean zero and variance one. Then, the expected value of the squared Frobenius norm of their product is given by $\EE({\norm{\Ab\Bb}^2_F}) = mnr$.
\begin{proof}
Using Proposition~\ref{prop:inner-expectation} we can write
\begin{align*}
 \EE({\norm{\Ab\Bb}^2_F})
 &= 
 \EE\left(\begin{tikzpicture}[baseline=\baseline]
   \tikzset{tensor/.style = {minimum size = 0.5cm,shape = circle,thick,draw=black,inner sep = 0pt}, edge/.style = {thick,line width=.4mm},every loop/.style={}}
    \def\y{0}
    \def\x{0}
    \node[tensor,fill=red!50!white] (A) at (\x,\y+0.5){\scalebox{0.85}{$\Ab$}};
    \node[tensor,fill=blue!50!white] (B) at (\x+1,\y+0.5){\scalebox{0.85}{$\Bb$}};
    \draw[edge] (A) -- (B) node [midway,above]{\scalebox{0.7}{\dimcolor{$r$}}};
    \node[tensor,fill=red!50!white] (A1) at (\x,\y-0.5){\scalebox{0.85}{$\Ab$}};
    \node[tensor,fill=blue!50!white] (B1) at (\x+1,\y-0.5){\scalebox{0.85}{$\Bb$}};
    \draw[edge] (A1) -- (B1) node [midway,above]{\scalebox{0.7}{\dimcolor{$r$}}};
    \draw[edge] (A) to [out=-150,in=160, distance=0.5cm] (A1);
    \draw[edge] (B) to [out=-30,in=20, distance=0.5cm] (B1);
    \node[draw=none] () at (\x-0.75,\y) {\scalebox{0.7}{\dimcolor{$m$}}};
    \node[draw=none] () at (\x+1.75,\y) {\scalebox{0.7}{\dimcolor{$n$}}};
    \end{tikzpicture}\right)\\
&=
\EE\left(\begin{tikzpicture}[baseline=\baseline]
   \tikzset{tensor/.style = {minimum size = 0.5cm,shape = circle,thick,draw=black,inner sep = 0pt}, edge/.style = {thick,line width=.4mm},every loop/.style={}}
    \def\y{0.25}
    \def\x{0}
    \node[tensor,fill=red!50!white] (A) at (\x,\y){\scalebox{0.85}{$\Ab$}}; 
    \node[tensor,fill=red!50!white] (A1) at (\x+1,\y){\scalebox{0.85}{$\Ab$}};
    \draw[edge](A)--(A1)node[midway,above] {\scalebox{0.7}{\dimcolor{$m$}}};
    \draw[edge](A)--(\x,\y-0.85)node[midway,left] {\scalebox{0.7}{\dimcolor{$r$}}};
    \draw[edge](A1)--(\x+1,\y-0.85)node[midway,right] {\scalebox{0.7}{\dimcolor{$r$}}};
    \node[tensor,fill=blue!50!white] (B) at (\x+2,\y){\scalebox{0.85}{$\Bb$}}; 
    \node[tensor,fill=blue!50!white] (B1) at (\x+3,\y){\scalebox{0.85}{$\Bb$}};
    \draw[edge](B)--(B1)node[midway,above] {\scalebox{0.7}{\dimcolor{$n$}}};
    \draw[edge](B)--(\x+2,\y-0.85)node[midway,left] {\scalebox{0.7}{\dimcolor{$r$}}};
    \draw[edge](B1)--(\x+3,\y-0.85)node[midway,right] {\scalebox{0.7}{\dimcolor{$r$}}};
    \draw[edge](\x+2,\y-0.85) -- (\x+2,\y-1) -- (\x+1,\y-1) -- (\x+1,\y-0.85);
    \draw[edge](\x+3,\y-0.85) -- (\x+3,\y-1.25) -- (\x,\y-1.25) -- (\x,\y-0.85);
    \draw[dotted](\x+1.5,\y+0.5) -- (\x+1.5,\y-1.65);
\end{tikzpicture}\right)
=
\begin{tikzpicture}[baseline=\baseline]
   \tikzset{tensor/.style = {minimum size = 0.5cm,shape = circle,thick,draw=black,inner sep = 0pt}, edge/.style = {thick,line width=.4mm},every loop/.style={}}
    \def\y{0}
    \def\x{0}
    \node[tensor,fill=red!50!white] (A) at (\x,\y){\scalebox{0.85}{$\Ab$}}; 
    \node[tensor,fill=red!50!white] (A1) at (\x+1,\y){\scalebox{0.85}{$\Ab$}};
    \draw[edge, draw=red](A)--(A1)node[midway,above] {\scalebox{0.7}{\dimcolor{$m$}}};
   \node[tensor,fill=blue!50!white] (B) at (\x+2,\y){\scalebox{0.85}{$\Bb$}}; 
    \node[tensor,fill=blue!50!white] (B1) at (\x+3,\y){\scalebox{0.85}{$\Bb$}};
    \draw[edge, draw=blue](B)--(B1)node[midway,above] {\scalebox{0.7}{\dimcolor{$n$}}};
    \node at (\x-0.55,\y) {$\EE\Biggl ( $};
    \node (eq) at (\x+1.39,\y) {$\Biggr )\EE$};
    \node (eq) at (\x+1.73,\y) {$\Biggl($};
    \draw[edge, draw=blue](B) -- (\x+2,\y-0.85);
    \draw[edge, draw=blue](B1) -- (\x+3,\y-0.85);
    \draw[edge, draw=red](A) -- (\x,\y-0.85);
    \draw[edge, draw=red](A1) -- (\x+1,\y-0.85);
    \draw[edge](\x+1,\y-0.85) -- (\x+1,\y-1) -- (\x+2,\y-1)--(\x+2,\y-0.85);
    \draw[edge](\x+3,\y-0.85) -- (\x+3,\y-1.2) -- (\x,\y-1.2) -- (\x,\y-0.85);
    \node[draw=none] () at (\x+1.65,\y-0.85) {\scalebox{0.7}{\dimcolor{$r$}}};
    \node[draw=none] () at (\x+0.5,\y-0.95) {\scalebox{0.7}{\dimcolor{$r$}}};
\end{tikzpicture}\Biggr).
\end{align*}
By Proposition~\ref{prop:wishart-mean}, we have $\EE\left(\begin{tikzpicture}[baseline=\baseline]
   \tikzset{tensor/.style = {minimum size = 0.5cm,shape = circle,thick,draw=black,inner sep = 0pt}, edge/.style = {thick,line width=.4mm},every loop/.style={}}
    \def\y{0.15}
    \def\x{0}
    \node[tensor,fill=red!50!white] (A) at (\x,\y){\scalebox{0.85}{$\Ab$}}; 
    \node[tensor,fill=red!50!white] (A1) at (\x+1,\y){\scalebox{0.85}{$\Ab$}};
    \draw[edge, draw=red](A)--(A1)node[midway,above] {\scalebox{0.7}{\dimcolor{$m$}}};
    \draw[edge, draw=red](A) -- (\x,\y-0.85)node[midway,left] {\scalebox{0.7}{\dimcolor{$r$}}};
    \draw[edge, draw=red](A1) -- (\x+1,\y-0.85)node[midway,right] {\scalebox{0.7}{\dimcolor{$r$}}};
\end{tikzpicture}\right)
=
m\begin{tikzpicture}[baseline=\baseline]
   \tikzset{tensor/.style = {minimum size = 0.5cm,shape = circle,thick,draw=black,inner sep = 0pt}, edge/.style = {thick,line width=.4mm},every loop/.style={}}
    \def\y{0}
    \def\x{0}
     \draw[edge, draw=red] (\x,\y-0.1) to [out=100,in=80] (\x+2,\y-0.1);
    \node[draw=none] () at (\x+1,\y+0.65) {\scalebox{0.7}{\dimcolor{$r$}}};
\end{tikzpicture}$
and 
$\EE\left(\begin{tikzpicture}[baseline=\baseline]
   \tikzset{tensor/.style = {minimum size = 0.5cm,shape = circle,thick,draw=black,inner sep = 0pt}, edge/.style = {thick,line width=.4mm},every loop/.style={}}
    \def\y{0.15}
    \def\x{0}
    \node[tensor,fill=blue!50!white] (B) at (\x,\y){\scalebox{0.85}{$\Bb$}}; 
    \node[tensor,fill=blue!50!white] (B1) at (\x+1,\y){\scalebox{0.85}{$\Bb$}};
    \draw[edge, draw=blue](B)--(B1)node[midway,above] {\scalebox{0.7}{\dimcolor{$n$}}};
    \draw[edge, draw=blue](B) -- (\x,\y-0.85)node[midway,left] {\scalebox{0.7}{\dimcolor{$r$}}};
    \draw[edge, draw=blue](B1) -- (\x+1,\y-0.85)node[midway,right] {\scalebox{0.7}{\dimcolor{$r$}}};
\end{tikzpicture}\right)
= 
n\begin{tikzpicture}[baseline=\baseline]
   \tikzset{tensor/.style = {minimum size = 0.5cm,shape = circle,thick,draw=black,inner sep = 0pt}, edge/.style = {thick,line width=.4mm},every loop/.style={}}
    \def\y{0}
    \def\x{0}
     \draw[edge, draw=blue] (\x,\y-0.1) to [out=100,in=80] (\x+2,\y-0.1);
    \node[draw=none] () at (\x+1,\y+0.65) {\scalebox{0.7}{\dimcolor{$r$}}};
\end{tikzpicture}.
$ 
Since there are contractions between the legs of size $r$ we finally obtain
\begin{align*}
 \EE({\norm{\Ab\Bb}^2_F})
 = 
 \EE\left(\begin{tikzpicture}[baseline=\baseline]
   \tikzset{tensor/.style = {minimum size = 0.5cm,shape = circle,thick,draw=black,inner sep = 0pt}, edge/.style = {thick,line width=.4mm},every loop/.style={}}
    \def\y{0}
    \def\x{0}
    \node[tensor,fill=red!50!white] (A) at (\x,\y+0.5){\scalebox{0.85}{$\Ab$}};
    \node[tensor,fill=blue!50!white] (B) at (\x+1,\y+0.5){\scalebox{0.85}{$\Bb$}};
    \draw[edge] (A) -- (B) node [midway,above]{\scalebox{0.7}{\dimcolor{$r$}}};
    \node[tensor,fill=red!50!white] (A1) at (\x,\y-0.5){\scalebox{0.85}{$\Ab$}};
    \node[tensor,fill=blue!50!white] (B1) at (\x+1,\y-0.5){\scalebox{0.85}{$\Bb$}};
    \draw[edge] (A1) -- (B1) node [midway,above]{\scalebox{0.7}{\dimcolor{$r$}}};
    \draw[edge] (A) to [out=-150,in=160, distance=0.5cm] (A1);
    \draw[edge] (B) to [out=-30,in=20, distance=0.5cm] (B1);
    \node[draw=none] () at (\x-0.75,\y) {\scalebox{0.7}{\dimcolor{$m$}}};
    \node[draw=none] () at (\x+1.75,\y) {\scalebox{0.7}{\dimcolor{$n$}}};
    \end{tikzpicture}\right)
=
mn\  \begin{tikzpicture}     [baseline=\baseline]
 \tikzset{tensor/.style = {minimum size = 0.5cm,shape = circle,thick,draw=black,inner sep = 0pt}, edge/.style = {thick,line 
width=.4mm},every loop/.style={}}
    \def\y{0}
    \def\x{1}
    \draw[edge](\x+2,\y) -- (\x+2,\y-0.35) -- (\x+3,\y-0.35)--(\x+3,\y);
    \draw[edge](\x+1,\y) -- (\x+1,\y-0.5) -- (\x+4,\y-0.5) -- (\x+4,\y);
    \draw[edge, draw=red] (\x+2,\y) to [out=100,in=80](\x+1,\y);
    \draw[edge, draw=blue] (\x+3,\y) to [out=100,in=80](\x+4,\y);
    \node[draw=none] () at (\x+2.5,\y-0.65) {\scalebox{0.7}{\dimcolor{$r$}}};
    \node[draw=none] () at (\x+3.5,\y+0.5) {\scalebox{0.7}{\dimcolor{$r$}}};
    \node[draw=none] () at (\x+1.5,\y+0.5) {\scalebox{0.7}{\dimcolor{$r$}}};
\end{tikzpicture}
= mnr.
\end{align*}
\end{proof}
\end{proposition}
Now that we understand how to compute the expected norm of products of Gaussian random matrices, we can extend this method to evaluate expectations of more complex tensor networks.

